% Осень26!A11; нумерация внутри столбца без пустых строк.
% Источник: лекция 4, §13 (строка 414)
\ticket{09}{Теорема о непрерывно дифференцируемых отображениях}
\label{ticket:09}

\subsection*{\texorpdfstring{Теорема о непрерывно дифференцируемых отображениях (класс \(C^1\))}{Теорема о непрерывно дифференцируемых отображениях (класс C\^{}1)}}

\textbf{Вводное.} Функциями многих переменных мы занимались и «учили новый язык»: в пространстве линейных операторов появилась \textbf{норма}, поэтому утверждения, где фигурирует производный оператор, удобно высказывать через норму.

\textbf{Напоминание (классы \(C^r\)).} \(F\in C^1(D)\): существуют все частные производные до первого порядка включительно во всех точках \(D\) и все они непрерывны на \(D\); аналогично \(C^r\) --- до порядка \(r\). Побочный продукт --- теорема о независимости частных производных от порядка дифференцирования.

\begin{theorem}

Пусть \(D\subset\mathbb{R}^m\) открыто, \(F:D\to\mathbb{R}^\ell\) дифференцируемо на \(D\) (координатные функции \(f_i\), \(i=1,\dots,\ell\)). Эквивалентны:

\begin{enumerate}
\def\labelenumi{\arabic{enumi}.}
\tightlist
\item
  \(F\in C^1(D)\): для всех \(i\in\overline{1,\ell}\), \(g\in\overline{1,m}\) существуют \(\dfrac{\partial f_i}{\partial x_g}(x)\) во всех точках \(x\in D\) и все они непрерывны на \(D\);
\item
  отображение \(F':D\to\mathrm{Lin}(\mathbb{R}^m,\mathbb{R}^\ell)\), \(F':x\mapsto F'(x)\) (точке сопоставляется производный оператор), \textbf{непрерывно} на \(D\) относительно операторной нормы: \[\forall a\in D\ \forall\varepsilon>0\ \exists\delta=\delta(\varepsilon,a)>0\ \forall x\in D:\ |x-a|<\delta\ \Rightarrow\ \|F'(x)-F'(a)\|<\varepsilon.\]
\end{enumerate}

Смысл: «продифференцировали отображение --- и получили непрерывный результат»; это и есть содержательное прочтение слов «непрерывно дифференцируемо».

\end{theorem}

\subsubsection*{Доказательство (1) ⇒ (2)}

Матрица оператора \(F'(x)-F'(a)\) (матрица Якоби разности) имеет элементы \(\dfrac{\partial f_i}{\partial x_g}(x)-\dfrac{\partial f_i}{\partial x_g}(a)\) --- всего \(m\ell\) штук. По лемме об оценке нормы оператора («прекрасное воспоминание»): \[\|F'(x)-F'(a)\|\le\sqrt{\sum_{i=1}^{\ell}\sum_{g=1}^{m}\Bigl(\frac{\partial f_i}{\partial x_g}(x)-\frac{\partial f_i}{\partial x_g}(a)\Bigr)^2}.\]

Пусть дано \(\varepsilon>0\). Каждая из \(m\ell\) частных производных непрерывна в точке \(a\): для каждой найдётся \(\delta_{i,g}>0\) такая, что \[|x-a|<\delta_{i,g}\ \Rightarrow\ \Bigl|\frac{\partial f_i}{\partial x_g}(x)-\frac{\partial f_i}{\partial x_g}(a)\Bigr|<\frac{\varepsilon}{\sqrt{m\ell}}.\] Положим \(\delta:=\min_{i,g}\delta_{i,g}\) (минимум конечного набора; «для каждой функции выбираю своё дельта и беру самое маленькое»). Тогда при \(|x-a|<\delta\) сразу все \(m\ell\) разностей по модулю \(<\varepsilon/\sqrt{m\ell}\), и \[\|F'(x)-F'(a)\|\le\sqrt{m\ell\cdot\frac{\varepsilon^2}{m\ell}}=\varepsilon.\qquad\blacksquare\]

\subsubsection*{Доказательство (2) ⇒ (1)}

Пусть дано \(\varepsilon>0\) и \(\delta\) из непрерывности \(F'\) в точке \(a\). \textbf{Прицелимся} вектором \(h=e_n\) (всюду нолики, на \(n\)-м месте единичка; \(|e_n|=1\)). Для \(|x-a|<\delta\): \[\bigl\|(F'(x)-F'(a))e_n\bigr\|\le\|F'(x)-F'(a)\|\cdot|e_n|<\varepsilon\] (общее свойство нормы оператора: \(\|Ah\|\le\|A\|\cdot|h|\)).

Но \((F'(x)-F'(a))e_n\) --- это \(n\)-й столбец матрицы, т.е. вектор в \(\mathbb{R}^\ell\) с координатами \[\frac{\partial f_i}{\partial x_n}(x)-\frac{\partial f_i}{\partial x_n}(a),\qquad i=1,\dots,\ell.\] Раз евклидова норма вектора (корень из суммы квадратов координат) меньше \(\varepsilon\), то и каждая координата по модулю меньше \(\varepsilon\): \[\Bigl|\frac{\partial f_i}{\partial x_n}(x)-\frac{\partial f_i}{\partial x_n}(a)\Bigr|<\varepsilon\qquad\forall i.\] Это \(\varepsilon\)-\(\delta\)-фраза непрерывности частной производной \(\dfrac{\partial f_i}{\partial x_n}\) в точке \(a\). В силу произвольности \(a\), \(n\), \(i\) --- все частные производные непрерывны на \(D\). \(\blacksquare\)

\emph{Мораль:} как и на первом курсе («для любого эпсилон --- дельта»), только добавилась размерность и контекст стали тяжелее; из \(a_1^2+a_2^2<\varepsilon^2\) следует \(|a_1|<\varepsilon\) --- здесь то же с \(m\ell\) слагаемыми.
